\documentclass[11pt,a4paper]{article}

\usepackage[utf8]{inputenc}
\usepackage[T1]{fontenc}
\usepackage[spanish,es-tabla,es-nodecimaldot]{babel}
\usepackage[margin=2.3cm]{geometry}
\usepackage{lmodern}
\usepackage{amsmath,amssymb}
\usepackage{booktabs,longtable,array,multirow}
\usepackage{xcolor}
\usepackage{graphicx}
\usepackage{tikz}
\usepackage{pgfplots}
\pgfplotsset{compat=1.17}
\usetikzlibrary{positioning,arrows.meta,shapes,calc,patterns}
\usepackage{listings}
\usepackage{enumitem}
\usepackage{caption}
\usepackage[hidelinks]{hyperref}
\usepackage{siunitx}
\usepackage{microtype}
\emergencystretch=2em
\newcommand{\pc}{\textcolor{red}{\textbf{?}}}
\sisetup{output-decimal-marker={,},group-separator={.},group-minimum-digits=4}

% ---------------------------------------------------------------------------
% Marcador de datos pendientes: todo lo que aparece en rojo debe completarse
% con mediciones reales antes de entregar.
\newcommand{\pend}[1]{\textcolor{red}{\textbf{[PENDIENTE: #1]}}}
% ---------------------------------------------------------------------------

\definecolor{hecho}{RGB}{46,125,50}
\definecolor{prog}{RGB}{230,140,0}
\definecolor{pendc}{RGB}{160,160,160}
\newcommand{\Hecho}{\textcolor{hecho}{Hecho}}
\newcommand{\EnProg}{\textcolor{prog}{En progreso}}
\newcommand{\Pend}{\textcolor{pendc}{Pendiente}}
\newcommand{\N}{$^{[N]}$}

\lstset{basicstyle=\ttfamily\small,columns=fullflexible,keepspaces=true,
  frame=single,framerule=0.3pt,rulecolor=\color{gray},xleftmargin=4pt,
  literate={á}{{\'a}}1 {é}{{\'e}}1 {í}{{\'i}}1 {ó}{{\'o}}1 {ú}{{\'u}}1 {ñ}{{\~n}}1}

\setlist{itemsep=2pt,topsep=4pt}

\title{Documentación trabajo de título ``forgethreads''\\[2pt]
\large Etapa 3: Mejoras, Pruebas Adicionales y Validación Preliminar}
\author{Carlos Riquelme Ávila}
\date{30 de octubre de 2026}

\begin{document}
\maketitle

\begin{abstract}
Este documento corresponde a la Etapa 3 del trabajo de título \textbf{forgethreads}, una
librería en CUDA C++ expuesta a Python para calcular efectos olvidados a gran escala en GPU.
Se presentan los ajustes derivados de la revisión de la Etapa 2, la secuencia de
optimizaciones aplicadas al kernel max--min (de la versión v1 a la v4 con \emph{tiling}
bidimensional), las pruebas de correctitud y de rendimiento realizadas, la validación
preliminar sobre matrices de \num{10000} nodos y la documentación técnica actualizada.
Al cierre de la etapa, el 100\% de las tareas de diseño y el 69\% de las de
implementación están completados.
\end{abstract}

% ===========================================================================
\section{Introducción}
% ===========================================================================
En las etapas anteriores se definió el problema (foRgotten~\cite{forgotten} y forgeffects no
alcanzan las dimensiones que exige la investigación del Dr.\ Julio Rojas Mora: \num{2640}
ejecuciones sobre matrices de hasta \num{10000} nodos), se diseñó la arquitectura en capas de
forgethreads y se implementó una primera versión correcta: el kernel v1, el orquestador
iterativo y los bindings PyBind11/DLPack.

La Etapa 3 tiene tres objetivos:
\begin{enumerate}
  \item \textbf{Optimizar} el cómputo partiendo de la línea base medida en la Etapa 2
        (kernel v1 limitado por L2), aplicando una mejora a la vez y midiéndola antes y
        después con Nsight Compute.
  \item \textbf{Probar} la solución: correctitud de cada variante frente a la referencia CPU,
        rendimiento frente a forgeffects y comportamiento en los tamaños objetivo.
  \item \textbf{Validar preliminarmente} el sistema completo en condiciones equivalentes a las
        de los experimentos de la Etapa 4 (servidor \texttt{cuda3}, RTX~4090, $n=\num{10000}$),
        y dejar preparada la entrega final.
\end{enumerate}

% ===========================================================================
\section{Incorporación de la retroalimentación}
% ===========================================================================
\pend{transcribir aquí los comentarios concretos del profesor guía sobre la Etapa 2 y, para
cada uno, el cambio aplicado. Los puntos siguientes recogen los compromisos que la propia
Etapa 2 dejó abiertos y cómo se resolvieron.}

\begin{itemize}
  \item \textbf{Optimización medida, no supuesta.} La Etapa 2 dejó como línea base el perfil
        del kernel v1 (Tabla~3 de esa etapa). En esta etapa cada mejora se aplicó por separado
        con el protocolo \emph{perfil antes $\rightarrow$ cambio $\rightarrow$ perfil después}
        documentado en \texttt{mejoras.md}, y se reportan también los resultados que no
        mejoraron lo esperado (Sección~\ref{sec:mejoras}).
  \item \textbf{Ajustes de diseño.} El orquestador pasó a trabajar con buffers
        \emph{padeados} a múltiplos del tamaño de tile y rellenos con un valor negativo
        (ADR-8), lo que permitió eliminar las guardas de borde del bucle interno de los kernels
        con \emph{tiling}. La interfaz pública no cambió: los bindings pasaron de invocar
        \texttt{maxmin} (v1) a \texttt{maxminv4} sin modificar la firma de
        \texttt{ft.maxmin}, como preveía la arquitectura en capas.
  \item \textbf{Modificaciones técnicas.} La reconstrucción de caminos (T2.1.2b) se integró en
        el orquestador y el resultado se entrega como la tupla
        \texttt{(paths, values, eta0)} prometida en los requerimientos transversales.
  \item \textbf{Planificación.} El backlog se actualizó con las tareas surgidas durante la
        etapa (marcadas con \N) y la carta Gantt de la Etapa 4 se rehízo con las
        desviaciones registradas (Sección~\ref{sec:plan}).
  \item \textbf{Repositorio.} Se adopta Conventional Commits~\cite{cc} con referencia al ID
        del backlog para todos los commits nuevos (Sección~\ref{sec:git}).
\end{itemize}

% ===========================================================================
\section{Registro de avance del backlog}
% ===========================================================================
La Tabla~\ref{tab:backlog} muestra el estado de cada tarea al cierre de la Etapa 3. Las
tareas marcadas con \N\ se agregaron durante esta etapa.

{\small
\begin{longtable}{p{1.45cm}p{7.3cm}p{1.6cm}p{1.75cm}p{0.9cm}}
\caption{Estado del backlog al cierre de la Etapa 3}\label{tab:backlog}\\
\toprule
\textbf{ID} & \textbf{Tarea} & \textbf{Tipo} & \textbf{Estado} & \textbf{Etapa}\\
\midrule\endfirsthead
\toprule
\textbf{ID} & \textbf{Tarea} & \textbf{Tipo} & \textbf{Estado} & \textbf{Etapa}\\
\midrule\endhead
\multicolumn{5}{l}{\emph{Épica 1 -- Límites de forgeffects (OBJ1)}}\\
T1.0.1\N & Diseñar la búsqueda de límites & Diseño & \Hecho & 1--2\\
T1.1.1 & Script de barrido de dimensión y threshold & Implem. & \EnProg & 4\\
T1.1.2 & Medición de tiempo con sincronización forzada & Implem. & \Hecho & 1--2\\
T1.1.3 & Ejecución planificada en Slurm sobre cuda3 & Implem. & \Hecho & 1--2\\
T1.1.4 & Exportación de resultados (media y desviación) & Implem. & \Hecho & 1--2\\
T1.1.5 & Perfilado de recursos de forgeffects & Validación & \Pend & 4\\
\midrule
\multicolumn{5}{l}{\emph{Épica 2 -- Diseño e implementación de forgethreads (OBJ2)}}\\
T2.0.1\N & Diseñar la arquitectura en capas & Diseño & \Hecho & 2\\
T2.1.1 & Estrategia de memoria y presupuesto de VRAM & Diseño & \Hecho & 2\\
T2.0.3\N & Diseñar el kernel v1 & Diseño & \Hecho & 2\\
T2.0.4\N & Diseñar la reconstrucción de caminos & Diseño & \Hecho & 2\\
T2.0.2\N & Contrato de interfaz y ciclo de vida DLPack & Diseño & \Hecho & 2\\
T2.0.5\N & Diseñar las variantes con \emph{tiling} (v2--v4) y el padding & Diseño & \Hecho & 3\\
T2.1.2a & Kernel v1 y orquestador iterativo & Implem. & \Hecho & 2\\
T2.1.4\N & Verificar el kernel v1 (CPU, compute-sanitizer) & Validación & \Hecho & 2\\
T2.3.1 & Bindings PyBind11 de \texttt{maxmin} & Implem. & \Hecho & 2\\
T2.3.3\N & Integración del módulo en CMake & Implem. & \Hecho & 2\\
T2.4.1 & Importación de tensores vía DLPack & Implem. & \Hecho & 2\\
T2.4.3\N & Exportación de buffers vía DLPack & Implem. & \Hecho & 2\\
T2.1.2b & Reconstrucción de caminos (host) & Implem. & \Hecho & 3\\
T2.5.1 & Empaquetado del resultado (caminos, valores, $\eta_0$) & Implem. & \Hecho & 3\\
T2.2.2 & Optimizar los kernels (shared memory, tiling) & Implem. & \Hecho & 3\\
T2.1.5\N & Verificar v4 frente a la referencia CPU & Validación & \Hecho & 3\\
T2.1.3 & Perfilar con Nsight y ajustar a los criterios de OBJ2 & Validación & \EnProg & 3--4\\
T2.2.1 & Benchmark comparativo frente a forgeffects & Implem. & \EnProg & 3--4\\
T2.1.2c\N & Reconstrucción de caminos en GPU (CSR de pivotes) & Implem. & \EnProg & 4\\
T2.4.2 & Prueba de ausencia de copias host--device & Validación & \Pend & 4\\
T2.3.2 & Empaquetado instalable vía pip & Implem. & \Pend & 4\\
T2.6.1 & Versiones mínimas de CUDA, driver y Python & Document. & \Pend & 4\\
\midrule
\multicolumn{5}{l}{\emph{Épica 3 -- Validación de la solución (OBJ3)}}\\
T3.1.1 & Diseñar el pipeline de experimentos de profundidad & Diseño & \Hecho & 2\\
T3.2.1 & Diseñar el protocolo de toma de tiempo & Diseño & \Hecho & 2\\
T3.2.2 & Implementar el protocolo de medición & Implem. & \EnProg & 3--4\\
T3.2.3 & Calcular y reportar el speedup & Validación & \EnProg & 3--4\\
T3.1.2 & Ejecutar y reportar las \num{2640} corridas & Validación & \Pend & 4\\
T3.3.1 & Análisis estadístico de los resultados & Validación & \Pend & 4\\
\midrule
\multicolumn{5}{l}{\emph{Documentación}}\\
T4.1\N & Documento de la Etapa 2 & Document. & \Hecho & 2\\
T4.3\N & Documento de la Etapa 3 & Document. & \Hecho & 3\\
T4.2\N & Informe final & Document. & \Pend & 4\\
\bottomrule
\end{longtable}
}

\begin{table}[h]
\centering
\caption{Avance por tipo de tarea (Etapa 2 $\rightarrow$ Etapa 3)}\label{tab:avance}
\begin{tabular}{lrrrrrr}
\toprule
\textbf{Tipo} & \textbf{Total} & \textbf{Hecho} & \textbf{En progreso} & \textbf{Pendiente} &
\textbf{\% Etapa 2} & \textbf{\% Etapa 3}\\
\midrule
Diseño         & 9  & 9  & 0 & 0  & 100\% & \textbf{100\%}\\
Implementación & 16 & 11 & 4 & 1  & 53\%  & \textbf{69\%}\\
Validación     & 9  & 2  & 2 & 5  & 14\%  & \textbf{22\%}\\
Documentación  & 4  & 2  & 0 & 2  & 33\%  & \textbf{50\%}\\
\midrule
Total          & 38 & 24 & 6 & 8  & 55\%  & \textbf{63\%}\\
\bottomrule
\end{tabular}
\end{table}

El avance de implementación (69\%) supera el 50\% exigido para esta etapa. Cada tarea
marcada como \emph{Hecho} tiene respaldo verificable en el repositorio
(Sección~\ref{sec:demo}).

% ===========================================================================
\section{Mejoras para optimizar el diseño}\label{sec:mejoras}
% ===========================================================================

\subsection{Punto de partida}
El perfil de Nsight Compute del kernel v1 en la Etapa 2 mostró una ocupancia de 99,16\%,
un throughput de cómputo de sólo 5,19\% y un throughput de L2 de 87,09\%: el kernel está
\textbf{limitado por el acceso a L2}. La causa es estructural: v1 asigna un bloque por celda
de salida, de modo que cada fila de $C_s$ y cada columna de $A$ se vuelven a pedir a la
jerarquía de memoria una vez por cada celda que las usa. El máximo y el mínimo cuestan una
instrucción cada uno; el kernel pasa casi todo su tiempo esperando datos.

La composición max--min tiene la misma estructura de accesos que un GEMM, con
$(\min,\max)$ en lugar de $(\times,+)$. Por eso la estrategia sigue la progresión de
Boehm~\cite{boehm} para GEMM: \emph{aumentar la reutilización de cada dato cargado}, primero
en shared memory y luego en registros, hasta que el kernel deje de estar limitado por la
memoria.

\subsection{Variantes implementadas}

\paragraph{v1b -- columna de $A$ en shared memory.}
Cada bloque copia una vez la columna $A[b,:,n]$ completa ($K$ elementos, 2\,KB para
$K=\num{1000}$) a shared memory mediante una carga cooperativa, y el bucle de reducción lee
de ahí. El bucle interno pasa de dos accesos globales por iteración a uno global y uno a
shared memory.

\paragraph{v2 -- tiling $32\times32$.}
Un bloque calcula un tile de $32\times32$ celdas de salida. En cada paso carga a shared memory
un tile de $C_s$ y uno de $A$, y cada hilo calcula su celda recorriendo los 32 elementos del
tile. Cada dato traído de memoria global se reutiliza 32 veces.

\paragraph{v3 -- \emph{blocktiling} 1D.}
Tiles de bloque de $64\times64$ ($BM=BN=64$, $BK=8$) y \textbf{8 resultados por hilo} en una
columna ($TM=8$). Cada hilo mantiene en registros los 8 máximos parciales con sus
\emph{argmax}; el valor de $A$ leído de shared memory se reutiliza 8 veces desde un registro.

\paragraph{v4 -- \emph{blocktiling} 2D.}
Mismos tiles de bloque, pero cada hilo calcula un sub-tile de $8\times8 = 64$ celdas
($TM=TN=8$), con 64 hilos por bloque. En cada paso del bucle interno el hilo carga 8 valores
de $C_s$ y 8 de $A$ desde shared memory a registros y evalúa los 64 pares
$\min(\cdot,\cdot)$ con su máximo y \emph{argmax}. La escritura final respeta los límites
lógicos $M\times N$ y aplica el umbral diferencial de la Ecuación~(2) de la Etapa 2.

\paragraph{Padding.}
Las variantes v2--v4 suponen dimensiones múltiplo del tile. El orquestador
\texttt{maxminv4} reserva los buffers con dimensiones \emph{padeadas} y los rellena con un
valor negativo (patrón \texttt{0xBC} en \texttt{half}). Como las incidencias están en
$[0,1]$, un valor negativo es neutro para el máximo y nunca supera el umbral, por lo que el
bucle interno no necesita condicionales de borde.

\subsection{Análisis de reutilización}
La Tabla~\ref{tab:reuso} cuantifica, para una celda de salida y un orden, cuántos accesos
requiere cada variante. Es un análisis independiente del hardware que explica por qué cada
paso debería reducir la presión sobre L2.

\begin{table}[h]
\centering
\caption{Accesos por celda de salida y por orden ($K$ = dimensión de la reducción)}
\label{tab:reuso}
\begin{tabular}{lccc}
\toprule
\textbf{Variante} & \textbf{Lecturas globales} & \textbf{Lecturas de shared} &
\textbf{Reutilización}\\
\midrule
v1   & $2K$        & 0           & 1\\
v1b  & $2K$        & $K$         & 1\\
v2 ($32\times32$)            & $2K/32$ & $2K$     & 32\\
v3 ($64\times64$, $TM=8$)    & $2K/64$ & $1{,}125\,K$ & 64\\
v4 ($64\times64$, $8\times8$) & $2K/64$ & $0{,}25\,K$  & 64\\
\bottomrule
\end{tabular}

\smallskip
{\footnotesize Reutilización: veces que se usa cada dato traído de memoria global. En v1b la
columna se lee de shared memory, lo que quita un acceso global del bucle interno.}
\end{table}

Para $n=\num{10000}$, v1 hace $2\cdot 10^4$ lecturas globales por celda y v4 hace 313; y v4
reduce además 8 veces las lecturas de shared memory respecto de v2, desplazando la
reutilización hacia los registros.

\subsection{Comparación antes y después}

\paragraph{Mejora v1 $\rightarrow$ v1b (medida).}
La Tabla~\ref{tab:v1b} resume el perfil de Nsight Compute antes y después del cambio, con
$B=1$, $n=\num{1000}$, $\theta=0{,}5$, en la RTX~4090 de \texttt{cuda3}. El conteo de
efectos coincide con la referencia CPU en ambos casos (\num{472326}).

\begin{table}[h]
\centering
\caption{Kernel v1 frente a v1b (Nsight Compute, $B=1$, $n=\num{1000}$)}\label{tab:v1b}
\begin{tabular}{lrrr}
\toprule
\textbf{Métrica} & \textbf{v1} & \textbf{v1b} & \textbf{$\Delta$}\\
\midrule
Duración [ms]                 & 14,58 & 11,39 & $-21{,}9\%$\\
Cómputo (SM) [\%]             & 5,19  & 10,45 & $+5{,}26$ pp\\
Memoria [\%]                  & 87,09 & 90,71 & $+3{,}62$ pp\\
L1/TEX [\%]                   & 37,87 & 39,64 & $+1{,}77$ pp\\
L2 [\%]                       & 87,09 & 90,71 & $+3{,}62$ pp\\
DRAM [\%]                     & 0,03  & 0,04  & $\approx$\\
Ocupancia lograda [\%]        & 99,16 & 99,55 & $\approx$\\
Registros por hilo            & 16    & 26    & $+10$\\
\bottomrule
\end{tabular}
\end{table}

La duración bajó un 22\%, pero \textbf{no por la razón esperada}. La hipótesis era que
cachear la columna reduciría la presión sobre L2; el throughput de L2 en cambio
\emph{subió}. Cada bloque sigue trayendo la columna completa una vez, así que el volumen total
de datos no cambia. La ganancia viene de quitar un acceso global del bucle interno, lo que
duplica el throughput de cómputo y hace que el kernel sature L2 durante menos tiempo. El
perfil de Nsight Systems de la misma configuración confirma un tiempo de kernel de 11,24\,ms.

Este resultado negativo parcial fue decisivo para el diseño: mostró que, mientras cada
bloque calcule una sola celda, ninguna caché por bloque reduce el volumen de datos. La
reutilización debe ocurrir \emph{entre celdas}, lo que justifica el cambio a \emph{tiling}
(v2--v4) en lugar de seguir ajustando v1.

\paragraph{Variantes v2--v4.}
\pend{completar las celdas marcadas con \pc{} en la Tabla~\ref{tab:variantes} con \texttt{scripts/profile.sh}
(\texttt{ncu --set full}) para cada variante en la misma configuración ($B=1$,
$n=\num{1000}$, $\theta=0{,}5$) y para $n=\num{10000}$.}

\begin{table}[h]
\centering
\caption{Perfil de las variantes (Nsight Compute, $B=1$, $n=\num{1000}$)}\label{tab:variantes}
\begin{tabular}{lrrrrr}
\toprule
\textbf{Métrica} & \textbf{v1} & \textbf{v1b} & \textbf{v2} & \textbf{v3} & \textbf{v4}\\
\midrule
Duración [ms]           & 14,58 & 11,39 & \pc & \pc & \pc\\
Cómputo (SM) [\%]       & 5,19  & 10,45 & \pc & \pc & \pc\\
L2 [\%]                 & 87,09 & 90,71 & \pc & \pc & \pc\\
Ocupancia lograda [\%]  & 99,16 & 99,55 & \pc & \pc & \pc\\
Speedup sobre v1        & 1,00$\times$ & 1,28$\times$ & \pc & \pc & \pc\\
\bottomrule
\end{tabular}
\end{table}

% ===========================================================================
\section{Avance de la implementación}\label{sec:demo}
% ===========================================================================

\subsection{Funcionalidades implementadas}
\begin{itemize}
  \item \textbf{Kernels} (\texttt{src/kernels/maxmin/}): \texttt{kernel\_v1.cu} (v1b),
        \texttt{kernel\_v2.cu}, \texttt{kernel\_v3\_multC.cu} y
        \texttt{kernel\_v4\_2D\_blocktiling.cu}, todos con la misma firma (valor, \emph{argmax}
        y contador), lo que permite cambiar de variante sin tocar los bindings.
  \item \textbf{Orquestadores} (\texttt{src/algorithms/maxmin\{,v2,v3,v4\}.cu}): bucle por
        orden con doble buffer, padding y convergencia anticipada ($\eta_0$).
  \item \textbf{Reconstrucción de caminos} (\texttt{src/algorithms/assemble\_paths.cu}): a
        partir del \emph{argmax} de cada orden, sin recalcular la composición.
  \item \textbf{Interfaz Python} (\texttt{src/bindings.cu}): \texttt{ft.maxmin} invoca
        \texttt{maxminv4} y devuelve \texttt{(paths, values, eta0)} como arreglos NumPy, un
        par por orden.
  \item \textbf{Reconstrucción en GPU} (rama \texttt{csr-paths-construction}, en progreso):
        kernels \texttt{count\_pivots} y \texttt{write\_pivots} que registran, por orden, los
        pivotes de cada celda en formato CSR, y un DAG de caminos (\texttt{path\_dag.cu}) que
        reconstruye los caminos de $m$ a $n$ sin copiar las matrices densas al host.
\end{itemize}

\subsection{Demostración}
\begin{lstlisting}[language=Python]
import numpy as np
import forgethreads as ft

a = np.load("input/CC.npy").astype(np.float16)[None]   # [1, n, n]
paths, values, eta0 = ft.maxmin(a, a, 0.5, 100)
print("profundidad de detención:", eta0)
for s, (p, v) in enumerate(zip(paths, values), start=1):
    print(f"orden {s}: {p.shape[0]} efectos, ancho de camino {p.shape[1]}")
\end{lstlisting}

El arnés C++ (\texttt{src/main.cu}) ejecuta la misma ruta con \texttt{maxminv4} y la compara
con la referencia CPU (Sección~\ref{sec:pruebas}).

% ===========================================================================
\section{Pruebas iniciales y resultados preliminares}\label{sec:pruebas}
% ===========================================================================

\subsection{Correctitud}
\begin{itemize}
  \item \textbf{v4 frente a la referencia CPU.} Con $B=1$, $n=\num{1000}$ y $\theta=0{,}5$,
        el arnés GPU (\texttt{src/main.cu}, que invoca \texttt{maxminv4}) encuentra
        \num{472326} efectos en el orden 1, exactamente el mismo conteo que la referencia CPU y
        que v1 (logs \texttt{fecuda.848} y \texttt{fecuda.850}). El cambio completo de
        estrategia de cómputo (bloque por celda $\rightarrow$ tiles $64\times64$ con padding)
        preserva el resultado. \pend{repetir la corrida registrando la variante en el log, ya
        que los logs actuales no la indican.}
  \item \textbf{Determinismo.} Ejecuciones repetidas sobre la misma entrada producen la misma
        secuencia de conteos por orden; por ejemplo, para $n=200$ y \texttt{order = 100} se
        obtiene siempre \num{2666}, \num{6502}, \num{7655}, \num{4000}, \ldots{} en las ocho
        corridas registradas con esa configuración (logs \texttt{fecuda.513} a \texttt{fecuda.926}).
  \item \textbf{compute-sanitizer.} memcheck, racecheck, initcheck y synccheck: 0 errores
        para v1 (Etapa 2). \pend{repetir los cuatro modos con v4 vía
        \texttt{scripts/sanitize.sh}.}
\end{itemize}

\subsection{Rendimiento preliminar frente a forgeffects}
Se midió el tiempo de extremo a extremo de \texttt{ft.maxmin} y de
\texttt{iterative\_maxmin\_cuadrado} de forgeffects (en GPU y en CPU) con \texttt{tests/test.py}:
$\theta=0{,}3$, 4 órdenes, 10 repeticiones tras calentamiento, matrices de $2n\times2n$ con
semilla fija, en \texttt{cuda3}. La Tabla~\ref{tab:prelim} y la Figura~\ref{fig:prelim}
muestran los resultados.

\begin{table}[h]
\centering
\caption{Tiempo de extremo a extremo (ms, media $\pm$ desviación estándar, 10 repeticiones)}
\label{tab:prelim}
\begin{tabular}{rrrr}
\toprule
\textbf{Nodos} & \textbf{forgeffects GPU} & \textbf{forgeffects CPU} & \textbf{ft.maxmin}\\
\midrule
200   & $28{,}1 \pm 6{,}0$     & $27{,}0 \pm 0{,}6$     & $5{,}4 \pm 0{,}7$\\
400   & $51{,}6 \pm 10{,}7$    & $75{,}8 \pm 2{,}2$     & $33{,}2 \pm 0{,}1$\\
600   & $107{,}7 \pm 28{,}5$   & $207{,}9 \pm 7{,}3$    & $111{,}7 \pm 0{,}8$\\
800   & $217{,}9 \pm 64{,}9$   & $430{,}3 \pm 9{,}4$    & $259{,}0 \pm 1{,}2$\\
1.000 & $394{,}3 \pm 111{,}4$  & $843{,}6 \pm 13{,}7$   & $497{,}5 \pm 1{,}6$\\
1.200 & $610{,}0 \pm 14{,}7$   & $\num{1457.9} \pm 48{,}4$  & $859{,}1 \pm 0{,}8$\\
1.600 & ---                    & $\num{3682.1} \pm 152{,}0$ & $\num{2039.3} \pm 5{,}2$\\
1.800 & ---                    & $\num{5158.1} \pm 220{,}7$ & $\num{2896.7} \pm 6{,}4$\\
2.000 & ---                    & ---                    & $\num{3969.3} \pm 7{,}8$\\
\bottomrule
\end{tabular}
\end{table}

\begin{figure}[h]
\centering
\begin{tikzpicture}
\begin{axis}[width=0.78\textwidth,height=6.3cm,xlabel={Nodos},ylabel={Tiempo [ms]},
  ymode=log,legend pos=south east,grid=both,grid style={gray!20},
  legend style={font=\small},xtick={200,600,1000,1400,1800},
  x tick label style={/pgf/number format/.cd,1000 sep={.}}]
\addplot+[mark=o,thick] coordinates {(200,28.1)(400,51.6)(600,107.7)(800,217.9)(1000,394.3)(1200,610.0)};
\addlegendentry{forgeffects GPU}
\addplot+[mark=triangle,thick] coordinates {(200,27.0)(400,75.8)(600,207.9)(800,430.3)(1000,843.6)(1200,1457.9)(1600,3682.1)(1800,5158.1)};
\addlegendentry{forgeffects CPU}
\addplot+[mark=square,thick] coordinates {(200,5.4)(400,33.2)(600,111.7)(800,259.0)(1000,497.5)(1200,859.1)(1600,2039.3)(1800,2896.7)(2000,3969.3)};
\addlegendentry{ft.maxmin}
\end{axis}
\end{tikzpicture}
\caption{Tiempo de extremo a extremo según el número de nodos (escala logarítmica)}
\label{fig:prelim}
\end{figure}

\paragraph{Análisis.}
\begin{itemize}
  \item forgethreads es 5,2 veces más rápido que forgeffects en GPU para 200 nodos, pero
        \textbf{a partir de 600 nodos es más lento}: 1,26 veces a \num{1000} nodos y 1,41 veces
        a \num{1200}. El objetivo de OBJ2 (speedup medio $\geq 10\times$) \textbf{no se cumple
        en esta medición}.
  \item El kernel no explica esa diferencia. Para $n=\num{1000}$ el kernel tarda 11,24\,ms por
        orden (Nsight Systems), unos 45\,ms en 4 órdenes, frente a 497\,ms de extremo a extremo.
        Más del 90\% del tiempo de \texttt{ft.maxmin} se gasta fuera del kernel. Los
        candidatos son la reconstrucción de caminos en el host, que copia $C_s$, $C_{s+1}$ y
        \emph{argmax} completos a la CPU en cada orden y recorre los caminos en un solo hilo,
        y la conversión de los caminos a listas de arreglos NumPy.
        \pend{confirmar el reparto con un perfil \texttt{nsys} de \texttt{ft.maxmin} completo.}
  \item forgethreads completa \num{1600}, \num{1800} y \num{2000} nodos, donde la medición de
        forgeffects en GPU ya no se ejecuta, y su desviación estándar es en general un orden de
        magnitud menor que la de forgeffects en GPU (0,1--1,6\,ms frente a 6--111\,ms), lo que
        indica tiempos mucho más estables.
\end{itemize}

\paragraph{Problemas de método detectados.}
La revisión del script reveló dos defectos que impiden usar esta tabla como el speedup oficial
(T3.2.3):
\begin{enumerate}
  \item \textbf{Entradas distintas.} forgeffects recibe matrices construidas con densidad
        $2n\log(2n)$ y forgethreads con $n\log n$. Ambas usan la misma semilla, pero no son la
        misma matriz, y la cantidad de efectos (y por lo tanto el costo de reconstruir
        caminos) difiere.
  \item \textbf{Versión del kernel no registrada.} El CSV no guarda el commit ni la variante
        de kernel usada, por lo que no es posible atribuir los tiempos a v1 o v4 con certeza.
\end{enumerate}

\subsection{Ajustes derivados de las pruebas}
\begin{itemize}
  \item \textbf{Reconstrucción de caminos en GPU} (T2.1.2c\N). Como el costo está fuera del
        kernel, optimizar más el kernel no cerraría la brecha. Se diseñó e implementó en la
        rama \texttt{csr-paths-construction} el registro de pivotes por orden en formato CSR
        directamente en GPU, de modo que sólo se transfieran al host los caminos que superan el
        umbral y no las matrices densas.
  \item \textbf{Dos modos de uso (planificado).} Los experimentos de profundidad de la Etapa 4
        sólo necesitan $\eta_0$ y los conteos por orden, no los caminos. Se agregará un modo que
        omita la reconstrucción, de modo que las \num{2640} corridas paguen sólo el costo del
        kernel (ADR-11).
  \item \textbf{Protocolo de medición corregido} (T3.2.2). La versión definitiva usa la misma
        matriz para ambas librerías, registra el commit y la variante en cada fila del CSV y
        separa tiempo de kernel (eventos CUDA) de tiempo de extremo a extremo.
\end{itemize}

% ===========================================================================
\section{Validación preliminar}\label{sec:validacion}
% ===========================================================================
La validación preliminar se ejecutó en el entorno real de la Etapa 4: servidor
\texttt{cuda3} (RTX~4090, 24\,GB), trabajos Slurm y matrices del tamaño objetivo.

\begin{table}[h]
\centering
\caption{Validación preliminar frente a los requerimientos}\label{tab:validacion}
\begin{tabular}{p{4.3cm}p{6.6cm}p{2.6cm}}
\toprule
\textbf{Requerimiento} & \textbf{Evidencia} & \textbf{Estado}\\
\midrule
Soportar $n=\num{10000}$ (crítico) &
Corridas completas con $n=\num{10000}$ y \texttt{order} entre 10 y 25 hasta la convergencia
(logs \texttt{fecuda.468}--\texttt{472}); por ejemplo, \num{26644}, \num{26711},
\num{21191}, \ldots{} efectos por orden. & Cumple\\
\addlinespace
Profundidad de detención en $n=\num{10000}$ &
$\eta_0 = 3$ con $c=50$ y $\eta_0 = 37$ con $c=2$
(\texttt{tests/test\_10000.csv}, rama de caminos). & Cumple\\
\addlinespace
Prueba de estrés de memoria &
Con $n=\num{100000}$ la guardia rechaza la reserva antes de lanzar el kernel
(``allocation for C\_dev\_after is \num{5640261632} bytes'', log \texttt{fecuda.473}): falla
de forma controlada, sin corromper el contexto CUDA. & Cumple\\
\addlinespace
Correctitud frente a CPU &
Conteo idéntico (\num{472326}) para v1 y v4. & Cumple\\
\addlinespace
Tiempo en $n=\num{10000}$ &
\num{26532}\,ms de extremo a extremo para 4 órdenes con $\theta=0{,}3$
(\texttt{output/timing\_results2.csv}, una corrida). & Preliminar\\
\addlinespace
Speedup medio $\geq 10\times$ (OBJ2) &
No se cumple de extremo a extremo (Tabla~\ref{tab:prelim}); en curso con la reconstrucción
en GPU. & No cumple aún\\
\addlinespace
Ocupancia y throughput $\geq 80\%$ (OBJ2) &
Ocupancia 99,55\% y memoria 90,71\% en v1b. \pend{medir v4.} & Parcial\\
\bottomrule
\end{tabular}
\end{table}

\paragraph{Estimación del costo de la Etapa 4.}
Con el tiempo medido en $n=\num{10000}$ (unos 6,6\,s por orden de extremo a extremo), las
corridas con profundidad alta (por ejemplo $\eta_0 = 37$ con $c = 2$) tardarían varios minutos
cada una. Sin reconstrucción de caminos, el costo por orden queda determinado por el kernel.
\pend{medir el tiempo del kernel v4 por orden en $n=\num{10000}$ y recalcular la duración
estimada de las \num{2640} corridas.}

% ===========================================================================
\section{Documentación del proyecto}
% ===========================================================================
\subsection{Decisiones técnicas (ADR)}
Se mantienen ADR-1 a ADR-6 de la Etapa 2 y se agregan:

{\small
\begin{longtable}{p{1.3cm}p{4.6cm}p{9.2cm}}
\caption{Nuevas decisiones de la Etapa 3}\\
\toprule
\textbf{ID} & \textbf{Decisión} & \textbf{Justificación y consecuencia}\\
\midrule
ADR-7 & \emph{Blocktiling} 2D (v4) como kernel por defecto &
Máxima reutilización por dato global y en registros (Tabla~\ref{tab:reuso}). Los bindings
cambian a \texttt{maxminv4} sin alterar la API. A cambio, el kernel es más complejo y requiere
dimensiones padeadas.\\
\addlinespace
ADR-8 & Padding a múltiplos del tile con relleno negativo &
Elimina las guardas de borde del bucle interno; el valor negativo es neutro para el máximo
porque las incidencias están en $[0,1]$. Cuesta hasta 63 filas y columnas extra de memoria
(48 para $n=\num{10000}$, menos del 1\%).\\
\addlinespace
ADR-9 & Mantener v1--v3 en el repositorio &
Sirven como referencia de correctitud y como evidencia de la progresión de optimización;
comparten firma con v4, por lo que pueden compararse en el mismo arnés.\\
\addlinespace
ADR-10 & Reconstrucción de caminos en GPU con CSR de pivotes &
Las pruebas mostraron que el costo de extremo a extremo está fuera del kernel. Registrar
pivotes en CSR en la GPU evita copiar matrices densas al host en cada orden.\\
\addlinespace
ADR-11 & Modo sin caminos para los experimentos &
Los experimentos de profundidad sólo necesitan $\eta_0$ y los conteos por orden; omitir los
caminos reduce el tiempo de las \num{2640} corridas. Se implementa en la Etapa 4.\\
\bottomrule
\end{longtable}
}

\subsection{Desviaciones del plan}
\begin{itemize}
  \item \textbf{Speedup oficial (T2.2.1, T3.2.3) desplazado a la Etapa 4.} La medición
        preliminar mostró que el cuello de botella está en la reconstrucción de caminos, y
        reveló defectos en el script. Reportar un speedup antes de corregirlos habría
        producido una cifra no comparable.
  \item \textbf{Nueva tarea T2.1.2c\N.} No estaba planificada; surgió del análisis de la
        Tabla~\ref{tab:prelim}.
  \item \textbf{Barrido de threshold y perfil de forgeffects (T1.1.1, T1.1.5) a la Etapa 4.}
        Se priorizó la optimización, que es el camino crítico para las \num{2640} corridas.
  \item \textbf{Prueba de ausencia de copias (T2.4.2) a la Etapa 4.} Hoy la salida de caminos
        pasa necesariamente por el host; la prueba sólo tiene sentido una vez integrada la
        reconstrucción en GPU.
\end{itemize}

\subsection{Repositorio git}\label{sec:git}
El proyecto se versiona en \url{https://github.com/criquelme2003/feCUDA}. El trabajo de esta
etapa se encuentra en las ramas \texttt{gpu\_paths} (kernels v2--v4 y orquestadores) y
\texttt{csr-paths-construction} (reconstrucción en GPU). Commits relevantes:
\texttt{cf44ae8} (kernel v2), \texttt{17f70a4} (kernels v3 y v4), \texttt{fa645cf}
(optimización de los kernels max--min), \texttt{322cef9} (corrección de caminos fallidos) y
\texttt{96c2da7} (registro CSR por orden y reconstrucción de caminos de $m$ a $n$).

Desde esta etapa los commits siguen Conventional Commits~\cite{cc}:
\begin{lstlisting}
perf(kernels): add 2D blocktiling maxmin kernel (v4)

Each thread computes an 8x8 sub-tile so values loaded from shared
memory are reused 8 times from registers. Matches CPU reference
(472326 effects, n=1000, thr=0.5).

Refs: T2.2.2
\end{lstlisting}

% ===========================================================================
\section{Planificación: Etapa 4}\label{sec:plan}
% ===========================================================================
\begin{figure}[h]
\centering
\begin{tikzpicture}[x=1.45cm,y=0.5cm,font=\small]
  % semanas: 0 = 2 nov, 1 = 9 nov, 2 = 16 nov, 3 = 23 nov, 4 = 30 nov, 5 = 4 dic
  \foreach \i/\l in {0/2 nov,1/9 nov,2/16 nov,3/23 nov,4/30 nov}
    \node[anchor=south] at (\i+0.5,0.3) {\l};
  \draw[gray!40] (0,0) -- (5,0);
  \foreach \i in {0,...,5} \draw[gray!30] (\i,0.3) -- (\i,-9.6);
  \def\gbar#1#2#3#4#5{\node[anchor=east] at (-0.1,-#1) {#5};
     \fill[#2] (#3,-#1-0.35) rectangle (#4,-#1+0.35);}
  \gbar{1}{blue!60}{0}{1.5}{T2.1.3 perfil Nsight de v2--v4}
  \gbar{2}{blue!60}{0}{2}{T2.1.2c caminos en GPU (integración)}
  \gbar{3}{blue!60}{0}{1}{T3.2.2 protocolo de medición corregido}
  \gbar{4}{blue!60}{1}{2}{T2.2.1 / T3.2.3 benchmark y speedup}
  \gbar{5}{blue!60}{1}{2}{T1.1.1 / T1.1.5 barrido y perfil FE}
  \gbar{6}{orange!80}{1}{3.5}{T3.1.2 \num{2640} corridas (Slurm)}
  \gbar{7}{orange!80}{2}{3}{T2.4.2, T2.3.2, T2.6.1 pip y README}
  \gbar{8}{orange!80}{3}{4}{T3.3.1 análisis estadístico}
  \gbar{9}{red!60}{3}{4.6}{T4.2 informe final}
  \fill[black] (4.6,-9.6) -- ++(0.12,0.25) -- ++(-0.24,0) -- cycle;
  \node[anchor=north] at (4.6,-9.7) {Entrega 4 dic};
\end{tikzpicture}
\caption{Carta Gantt de la Etapa 4}\label{fig:gantt}
\end{figure}

La Etapa 4 empieza por cerrar lo que la Etapa 3 dejó abierto: perfilar v2--v4, integrar la
reconstrucción de caminos en GPU y medir el speedup con el protocolo corregido. Con eso se
lanzan las \num{2640} corridas en Slurm, en paralelo con el empaquetado. El análisis
estadístico y el informe final cierran el trabajo antes del 4 de diciembre.

\begin{thebibliography}{9}
\bibitem{kaufmann} A. Kaufmann y J. Gil Aluja, \emph{Modelos para la investigación de efectos
  olvidados}. Vigo: Milladoiro, 1988.
\bibitem{forgotten} E. Mardones-Arias y J. Rojas-Mora, \emph{foRgotten: paquete de GNU R para
  el análisis de efectos olvidados}, 2022.
\bibitem{boehm} S. Boehm, «How to Optimize a CUDA Matmul Kernel for cuBLAS-like Performance:
  a Worklog», 2022. [En línea]. Disponible: \url{https://siboehm.com/articles/22/CUDA-MMM}
\bibitem{bpg} NVIDIA Corporation, \emph{CUDA C++ Best Practices Guide}, v12.5, 2024.
\bibitem{ncu} NVIDIA Corporation, \emph{Nsight Compute Documentation}, 2024. [En línea].
  Disponible: \url{https://docs.nvidia.com/nsight-compute/}
\bibitem{nygard} M. Nygard, «Documenting Architecture Decisions», 2011. [En línea].
  Disponible: \url{https://cognitect.com/blog/2011/11/15/documenting-architecture-decisions}
\bibitem{cc} «Conventional Commits 1.0.0». [En línea]. Disponible:
  \url{https://www.conventionalcommits.org}
\end{thebibliography}

\end{document}
